\documentclass[10pt]{exam}
\usepackage[hon]{template-for-exam}
\usepackage{silence}
\WarningFilter{latex}{Label `question}
\WarningFilter{latex}{There were multiply-defined labels}


\title{Magnetism Stations}
\author{Rohrbach}
\date{\today}

\begin{document}
\maketitle


\section{Magnetic Fields}

\begin{questions}
  \question \textbf{Bar Magnet:} Place a bar magnet under a piece of paper and sprinkle some iron filings on the paper.  \emph{Please do not let the iron filing touch the magnet because you will not be able to get them off!}
  \begin{parts}
    \part Draw the pattern you see on your paper in the space below. Make sure to label the North and South poles.
    \part Use the compass to draw the direction of the magnetic field. The North arrow on the compass points in the direction of the field.
  \end{parts}
  \vs[3]


\question \textbf{Horseshoe Magnet:} Carefully lift the paper up and remove the bar magnet. Empty the iron filings back into your beaker.  Now, place a new blank sheet of paper over the top of your horseshoe magnet. 
  \begin{parts}
    \part Draw the pattern you see on your paper in the space below. Make sure to label the North and South poles.
    \part Use the compass to draw the direction of the magnetic field. The North arrow on the compass points in the direction of the field.
    \part Make sure to look at the field both on the inside and outside of the horseshoe magnet.
  \end{parts}

  \vs[3]

\question Given your drawings, what can you say about the direction of the magnetic field?  How are magnetic field lines similar to and different than electric field lines?

  \vs



  \pagebreak



\question \textbf{Repelling Magnets:} Carefully lift the paper up, remove the button magnet, and place one of your two papers on top of two bar magnets, parallel to one another, separated by a spacer (the plastic piece that was between them in the box), and with their North poles in the same direction.

  \begin{parts}
    \part Draw the pattern you see on your paper in the space below
    \part Use the compass to draw the direction of the magnetic field.
  \end{parts}

  \vs[3]

\question \textbf{Attracting Magnets:} Carefully lift the paper up, change the orientation of one of your magnets (so the North poles point in opposite directions), and place your other paper on the magnets. Gently shake the filings onto the paper covering the magnets

  \begin{parts}
    \part Draw the pattern you see on your paper in the space below
    \part Use the compass to draw the direction of the magnetic field.
  \end{parts}


  \vs[3]


\end{questions}


\pagebreak

\section{Electromagnet}

\begin{questions}
\question Read about Solenoids and Electromagnets on pp. 637-638.  Define these terms:

\vspace{2em}

solenoid: \vs

electromagnet: \vs 


\question What three things affect the strength of an electromagnet (make sure to read the top of p. 638 to answer this)? \vs 

\question Try wrapping the wire around a nail 10 times as shown Try to pick up the paper clips.  How many paper clips were you able to pick up?  

\vspace{4em}

\question Using what you read in the book, predict what will happen if you wrap more wire around the nail.

\vs 

\question Wrap the wire around the nail 20 times.  How many paper clips were you able to pick up? 

\vspace{4em}


\question Was your prediction correct? Write down what you observe.

\vs

\vspace{5em}


\end{questions}


\pagebreak 



\section{Induced Voltage}


Using a wire, connect the red terminal on the voltmeter to the red, 300T terminal on the solenoid. Next, connect the black terminal on the voltmeter to the black 0 terminal on the solenoid. The measurements the voltmeter gives you will be in millivolts.  In each of the following experiments, pay attention to (a) the maximum number of millivolts you can induce and (b) the sign of the current [positive vs. negative].  Record your observations below.

\begin{questions}

  \question Experiment with inserting the magnet quickly versus slowly into the coil.  Which gives you a greater voltage? \vs 

  \question Experiment with putting the north pole of the magnet and the south pole.  What does this appear to do to the voltage? \vs

  \question Insert your magnet in the coil and don't move it.  What happens? \vs 

  \question Experiment with inserting the magnet on the two opposite sides of the solenoid.  What happens? \vs 

  \question Experiment with moving the magnet over the top of the solenoid, as opposed to through it.  Does this generate a voltage?  Is it more or less than when you move the magnet through the inside of the solenoid?  \vs 

  \question Change your wire to start at 100T, then change to 300T and eventually 500T.  This corresponds to the number of turns in the solenoid that are being used.  What happens as you increase the number of turns? \vs

\end{questions} 


\end{document}

